\section{Where parsing actually happens}
\label{sec:ch03-parsing}
\index{parsing|(}

Number five in that list is \code{DocumentParserMiddleware}. It is the obvious
place to look for the moment your query stops being text and starts being a
syntax tree, and for a GraphQL server reached over HTTP it is the wrong place.

Here is the measurement that gave it away. Every timeline Mosaic has printed,
against a cold cache or a warm one, reports the parse phase as a dash:

\begin{minted}[fontsize=\footnotesize]{text}
parse - validate 0.204ms compile 0.097ms coerce - execute 0.492ms total 0.931ms
    (document cache miss, operation cache miss, 146 resolvers)
\end{minted}

A dash means the timer never started. That first line is a document the server
had never seen, so nothing could have been cached, and the parse phase still did
not run. Something else had already done it.

\index{parsing!in the transport}%
That something is the transport. HotChocolate's ASP.NET Core layer reads the
request body with a \code{Utf8GraphQLRequestParser}, handing it the document
cache and the hash provider, and inside that parser the document is turned into
a tree before the executor is involved at all:

\begin{minted}{csharp}
var hash = _hashProvider!.ComputeHash(documentBody);
...
document = _cache!.TryGetDocument(hash.Value, out cachedDocument)
    ? cachedDocument.Body
    : Utf8GraphQLParser.Parse(documentBody, _options);
\end{minted}

By the time the pipeline runs, the request is carrying a parsed document rather
than source text. \code{DocumentParserMiddleware} checks for exactly that, and
its diagnostic scope lives inside the other branch:

\begin{minted}{csharp}
if (query is OperationDocument parsed)
{
    documentInfo.Hash = CreateDocumentHash(documentInfo, context.Request);
    documentInfo.Document = parsed.Document;
    ...
}
else if (query is OperationDocumentSourceText source)
{
    using (_diagnosticEvents.ParseDocument(context))
    {
        ...
    }
}
\end{minted}

So the middleware is not dead code and the phase is not fake. Both exist for
callers that hand the executor raw source text, which is what you get when you
execute a request in process rather than over HTTP, and
chapter~\ref{ch:automated-tests} does exactly that. Over HTTP the branch is
simply never taken.

\subsection*{A syntax error never gets in}

\index{parsing!syntax errors}%
The consequence is sharper than a missing timer, and it is testable. Send
Mosaic something that does not parse:

\begin{minted}{graphql}
{ products { title
\end{minted}

\begin{minted}[fontsize=\footnotesize]{json}
{
  "errors": [
    {
      "message": "Expected a `RightBrace`-token, but found a `EndOfFile`-token.",
      "locations": [ { "line": 1, "column": 20 } ],
      "extensions": { "code": "HC0011" }
    }
  ]
}
\end{minted}

HTTP 400, error \code{HC0011}, and no timeline line at all. I checked that
rather than inferring it: one good request followed by three different
unparseable documents against a warmed-up service produced exactly one timeline
line, from the good request. The three failures never reached the execution
pipeline, so the instrumentation middleware at position one never opened a
scope, so nothing was there to report.

\index{errors!layering}%
Put that next to the other rejection chapter~\ref{ch:hotchocolate-quickly}
showed, the field that does not exist:

\begin{minted}[fontsize=\footnotesize]{json}
{
  "errors": [
    {
      "message": "The field `nope` does not exist on the type `Product`.",
      "locations": [ { "line": 1, "column": 14 } ],
      "extensions": {
        "type": "Product",
        "field": "nope",
        "responseName": "nope",
        "specifiedBy": "https://spec.graphql.org/September2025/#sec-Field-Selections"
      }
    }
  ]
}
\end{minted}

Both are 400. Both have no \code{data} key. One carries a HotChocolate error
code and no \code{specifiedBy}; the other carries \code{specifiedBy} and no
code. Chapter~\ref{ch:hotchocolate-quickly} read that difference as the language
rules against the server's own, which is right, and I described both as
validation running before any resolver, which is not quite. The unknown field is
validation, at position six, inside the pipeline. The syntax error is the
transport refusing to build a request at all.

The correction is small and the reason to make it is not. From
chapter~\ref{ch:enter-the-router} onward there is a router in front of these
services, and \enquote{which component produced this error} becomes a question
you ask several times a week. A server that answers 400 without ever creating a
request is a server whose logs, traces and diagnostic listeners will all be
empty for that request. If you are looking for it in the wrong layer you will
conclude the request never arrived, and you will be half right in the way that
wastes the most time.

\index{parsing!hashing}%
One more thing falls out of the parser code above, and the next section is
built on it. Before the transport parses anything it computes a hash of the
document body and looks that hash up. The document cache is therefore consulted
twice per request, once at the door and once inside, and the second consultation
is the one that decides whether your query gets validated.
\index{parsing|)}
